% LaTeX Figure Templates for SHAKTI-CHAIN Publication
% Include with: % LaTeX Figure Templates for SHAKTI-CHAIN Publication
% Include with: % LaTeX Figure Templates for SHAKTI-CHAIN Publication
% Include with: % LaTeX Figure Templates for SHAKTI-CHAIN Publication
% Include with: \input{figures.tex}

% Required packages (should be in main document)
% \usepackage{graphicx}
% \usepackage{subcaption}
% \usepackage{tikz}
% \usepackage{pgfplots}

% ============================================================
% FIGURE PATHS
% ============================================================
\graphicspath{{figures/}{./figures/}{../figures/}}

% ============================================================
% SINGLE COLUMN FIGURE
% ============================================================
% Usage: \singlefig{filename}{caption}{label}

\newcommand{\singlefig}[3]{
\begin{figure}[htbp]
    \centering
    \includegraphics[width=\columnwidth]{#1}
    \caption{#2}
    \label{#3}
\end{figure}
}

% ============================================================
% DOUBLE COLUMN FIGURE
% ============================================================
% Usage: \doublefig{filename}{caption}{label}

\newcommand{\doublefig}[3]{
\begin{figure*}[htbp]
    \centering
    \includegraphics[width=\textwidth]{#1}
    \caption{#2}
    \label{#3}
\end{figure*}
}

% ============================================================
% SIDE-BY-SIDE FIGURES
% ============================================================
% Usage: \sidebyside{file1}{cap1}{lab1}{file2}{cap2}{lab2}{main_caption}{main_label}

\newcommand{\sidebyside}[8]{
\begin{figure}[htbp]
    \centering
    \begin{subfigure}[b]{0.48\columnwidth}
        \centering
        \includegraphics[width=\textwidth]{#1}
        \caption{#2}
        \label{#3}
    \end{subfigure}
    \hfill
    \begin{subfigure}[b]{0.48\columnwidth}
        \centering
        \includegraphics[width=\textwidth]{#4}
        \caption{#5}
        \label{#6}
    \end{subfigure}
    \caption{#7}
    \label{#8}
\end{figure}
}

% ============================================================
% THREE PANEL FIGURE
% ============================================================
% Usage: \threepanel{f1}{c1}{l1}{f2}{c2}{l2}{f3}{c3}{l3}{main_cap}{main_lab}

\newcommand{\threepanel}[2]{
\begin{figure*}[htbp]
    \centering
    #1
    \caption{#2}
\end{figure*}
}

% ============================================================
% HYPOTHESIS SUMMARY BAR CHART (TikZ)
% ============================================================
% Requires: \usepackage{pgfplots}
% Usage: Include data directly in document

% Example:
% \begin{figure}[htbp]
%     \centering
%     \begin{tikzpicture}
%         \begin{axis}[
%             ybar stacked,
%             bar width=15pt,
%             symbolic x coords={Domain1, Domain2, Domain3},
%             xtick=data,
%             xlabel={Research Domain},
%             ylabel={Number of Hypotheses},
%             legend style={at={(0.5,-0.15)}, anchor=north},
%         ]
%             \addplot[fill=shaktigreen] coordinates {(Domain1,8) (Domain2,9) (Domain3,7)};
%             \addplot[fill=shaktired] coordinates {(Domain1,2) (Domain2,3) (Domain3,3)};
%             \legend{Supported, Not Supported}
%         \end{axis}
%     \end{tikzpicture}
%     \caption{Summary of hypothesis test results by domain.}
%     \label{fig:hypothesis_summary}
% \end{figure}

% ============================================================
% BOX PLOT TEMPLATE
% ============================================================
% For efficiency comparisons - use with pgfplots

% Example:
% \begin{figure}[htbp]
%     \centering
%     \begin{tikzpicture}
%         \begin{axis}[
%             boxplot/draw direction=y,
%             xtick={1,2,3,4},
%             xticklabels={SHAKTI, Fixed, Uniform, CDA},
%             ylabel={Efficiency},
%         ]
%             \addplot+[boxplot prepared={
%                 median=0.94,
%                 upper quartile=0.96,
%                 lower quartile=0.91,
%                 upper whisker=0.99,
%                 lower whisker=0.85
%             }] coordinates {};
%             % Add more systems...
%         \end{axis}
%     \end{tikzpicture}
%     \caption{Efficiency comparison across trading mechanisms.}
%     \label{fig:efficiency_box}
% \end{figure}

% ============================================================
% SCALABILITY PLOT (LOG-LOG)
% ============================================================
% For showing O(n log n) scaling behavior

% Example:
% \begin{figure}[htbp]
%     \centering
%     \begin{tikzpicture}
%         \begin{loglogaxis}[
%             xlabel={Number of Agents},
%             ylabel={Clearing Time (ms)},
%             legend pos=north west,
%         ]
%             \addplot[only marks, mark=*] table {data/scalability.dat};
%             \addplot[domain=10:10000, samples=100, red, dashed] {x*ln(x)/100};
%             \legend{Observed, $O(n \log n)$}
%         \end{loglogaxis}
%     \end{tikzpicture}
%     \caption{Scalability analysis showing clearing time vs number of agents.}
%     \label{fig:scalability}
% \end{figure}

% ============================================================
% PARETO FRONT PLOT
% ============================================================
% For showing efficiency vs fairness trade-off

% Example:
% \begin{figure}[htbp]
%     \centering
%     \begin{tikzpicture}
%         \begin{axis}[
%             xlabel={Efficiency},
%             ylabel={Fairness (Gini Index)},
%             xmin=0.6, xmax=1.0,
%             ymin=0.6, ymax=1.0,
%         ]
%             \addplot[only marks, mark=o] coordinates {
%                 (0.78, 0.65) (0.82, 0.72) (0.89, 0.85) (0.94, 0.91)
%             };
%             \addplot[thick, shaktiblue, mark=*] coordinates {
%                 (0.94, 0.91)
%             };
%             \node[anchor=west] at (axis cs:0.94,0.91) {SHAKTI};
%         \end{axis}
%     \end{tikzpicture}
%     \caption{Pareto front showing efficiency-fairness trade-off.}
%     \label{fig:pareto}
% \end{figure}

% ============================================================
% VIOLIN PLOT APPROXIMATION
% ============================================================
% TikZ doesn't have native violin plots, use matplotlib export

% ============================================================
% FIGURE PLACEMENT HINTS
% ============================================================
% [h] - here (approximately)
% [t] - top of page
% [b] - bottom of page
% [p] - separate page for floats
% [!] - override internal parameters
%
% Recommended: [htbp] for flexibility

% ============================================================
% FIGURE SIZING GUIDELINES
% ============================================================
% Single column: \columnwidth (typically 3.5 inches)
% Double column: \textwidth (typically 7 inches)
% Specific sizes:
%   - Quarter page: 0.45\columnwidth
%   - Half page: 0.9\columnwidth
%   - Full width: \textwidth

% ============================================================
% RESOLUTION REQUIREMENTS
% ============================================================
% IEEE requires:
%   - Line art: 600 dpi minimum
%   - Halftones: 300 dpi minimum
%   - Color figures: 300 dpi minimum, CMYK color space
%
% Export settings:
%   - PDF for vector graphics (preferred)
%   - PNG at 300+ dpi for raster
%   - Avoid JPEG for figures with text/lines

% ============================================================
% ACCESSIBILITY
% ============================================================
% Ensure figures are interpretable in grayscale
% Use patterns in addition to colors
% Provide detailed captions describing key findings


% Required packages (should be in main document)
% \usepackage{graphicx}
% \usepackage{subcaption}
% \usepackage{tikz}
% \usepackage{pgfplots}

% ============================================================
% FIGURE PATHS
% ============================================================
\graphicspath{{figures/}{./figures/}{../figures/}}

% ============================================================
% SINGLE COLUMN FIGURE
% ============================================================
% Usage: \singlefig{filename}{caption}{label}

\newcommand{\singlefig}[3]{
\begin{figure}[htbp]
    \centering
    \includegraphics[width=\columnwidth]{#1}
    \caption{#2}
    \label{#3}
\end{figure}
}

% ============================================================
% DOUBLE COLUMN FIGURE
% ============================================================
% Usage: \doublefig{filename}{caption}{label}

\newcommand{\doublefig}[3]{
\begin{figure*}[htbp]
    \centering
    \includegraphics[width=\textwidth]{#1}
    \caption{#2}
    \label{#3}
\end{figure*}
}

% ============================================================
% SIDE-BY-SIDE FIGURES
% ============================================================
% Usage: \sidebyside{file1}{cap1}{lab1}{file2}{cap2}{lab2}{main_caption}{main_label}

\newcommand{\sidebyside}[8]{
\begin{figure}[htbp]
    \centering
    \begin{subfigure}[b]{0.48\columnwidth}
        \centering
        \includegraphics[width=\textwidth]{#1}
        \caption{#2}
        \label{#3}
    \end{subfigure}
    \hfill
    \begin{subfigure}[b]{0.48\columnwidth}
        \centering
        \includegraphics[width=\textwidth]{#4}
        \caption{#5}
        \label{#6}
    \end{subfigure}
    \caption{#7}
    \label{#8}
\end{figure}
}

% ============================================================
% THREE PANEL FIGURE
% ============================================================
% Usage: \threepanel{f1}{c1}{l1}{f2}{c2}{l2}{f3}{c3}{l3}{main_cap}{main_lab}

\newcommand{\threepanel}[2]{
\begin{figure*}[htbp]
    \centering
    #1
    \caption{#2}
\end{figure*}
}

% ============================================================
% HYPOTHESIS SUMMARY BAR CHART (TikZ)
% ============================================================
% Requires: \usepackage{pgfplots}
% Usage: Include data directly in document

% Example:
% \begin{figure}[htbp]
%     \centering
%     \begin{tikzpicture}
%         \begin{axis}[
%             ybar stacked,
%             bar width=15pt,
%             symbolic x coords={Domain1, Domain2, Domain3},
%             xtick=data,
%             xlabel={Research Domain},
%             ylabel={Number of Hypotheses},
%             legend style={at={(0.5,-0.15)}, anchor=north},
%         ]
%             \addplot[fill=shaktigreen] coordinates {(Domain1,8) (Domain2,9) (Domain3,7)};
%             \addplot[fill=shaktired] coordinates {(Domain1,2) (Domain2,3) (Domain3,3)};
%             \legend{Supported, Not Supported}
%         \end{axis}
%     \end{tikzpicture}
%     \caption{Summary of hypothesis test results by domain.}
%     \label{fig:hypothesis_summary}
% \end{figure}

% ============================================================
% BOX PLOT TEMPLATE
% ============================================================
% For efficiency comparisons - use with pgfplots

% Example:
% \begin{figure}[htbp]
%     \centering
%     \begin{tikzpicture}
%         \begin{axis}[
%             boxplot/draw direction=y,
%             xtick={1,2,3,4},
%             xticklabels={SHAKTI, Fixed, Uniform, CDA},
%             ylabel={Efficiency},
%         ]
%             \addplot+[boxplot prepared={
%                 median=0.94,
%                 upper quartile=0.96,
%                 lower quartile=0.91,
%                 upper whisker=0.99,
%                 lower whisker=0.85
%             }] coordinates {};
%             % Add more systems...
%         \end{axis}
%     \end{tikzpicture}
%     \caption{Efficiency comparison across trading mechanisms.}
%     \label{fig:efficiency_box}
% \end{figure}

% ============================================================
% SCALABILITY PLOT (LOG-LOG)
% ============================================================
% For showing O(n log n) scaling behavior

% Example:
% \begin{figure}[htbp]
%     \centering
%     \begin{tikzpicture}
%         \begin{loglogaxis}[
%             xlabel={Number of Agents},
%             ylabel={Clearing Time (ms)},
%             legend pos=north west,
%         ]
%             \addplot[only marks, mark=*] table {data/scalability.dat};
%             \addplot[domain=10:10000, samples=100, red, dashed] {x*ln(x)/100};
%             \legend{Observed, $O(n \log n)$}
%         \end{loglogaxis}
%     \end{tikzpicture}
%     \caption{Scalability analysis showing clearing time vs number of agents.}
%     \label{fig:scalability}
% \end{figure}

% ============================================================
% PARETO FRONT PLOT
% ============================================================
% For showing efficiency vs fairness trade-off

% Example:
% \begin{figure}[htbp]
%     \centering
%     \begin{tikzpicture}
%         \begin{axis}[
%             xlabel={Efficiency},
%             ylabel={Fairness (Gini Index)},
%             xmin=0.6, xmax=1.0,
%             ymin=0.6, ymax=1.0,
%         ]
%             \addplot[only marks, mark=o] coordinates {
%                 (0.78, 0.65) (0.82, 0.72) (0.89, 0.85) (0.94, 0.91)
%             };
%             \addplot[thick, shaktiblue, mark=*] coordinates {
%                 (0.94, 0.91)
%             };
%             \node[anchor=west] at (axis cs:0.94,0.91) {SHAKTI};
%         \end{axis}
%     \end{tikzpicture}
%     \caption{Pareto front showing efficiency-fairness trade-off.}
%     \label{fig:pareto}
% \end{figure}

% ============================================================
% VIOLIN PLOT APPROXIMATION
% ============================================================
% TikZ doesn't have native violin plots, use matplotlib export

% ============================================================
% FIGURE PLACEMENT HINTS
% ============================================================
% [h] - here (approximately)
% [t] - top of page
% [b] - bottom of page
% [p] - separate page for floats
% [!] - override internal parameters
%
% Recommended: [htbp] for flexibility

% ============================================================
% FIGURE SIZING GUIDELINES
% ============================================================
% Single column: \columnwidth (typically 3.5 inches)
% Double column: \textwidth (typically 7 inches)
% Specific sizes:
%   - Quarter page: 0.45\columnwidth
%   - Half page: 0.9\columnwidth
%   - Full width: \textwidth

% ============================================================
% RESOLUTION REQUIREMENTS
% ============================================================
% IEEE requires:
%   - Line art: 600 dpi minimum
%   - Halftones: 300 dpi minimum
%   - Color figures: 300 dpi minimum, CMYK color space
%
% Export settings:
%   - PDF for vector graphics (preferred)
%   - PNG at 300+ dpi for raster
%   - Avoid JPEG for figures with text/lines

% ============================================================
% ACCESSIBILITY
% ============================================================
% Ensure figures are interpretable in grayscale
% Use patterns in addition to colors
% Provide detailed captions describing key findings


% Required packages (should be in main document)
% \usepackage{graphicx}
% \usepackage{subcaption}
% \usepackage{tikz}
% \usepackage{pgfplots}

% ============================================================
% FIGURE PATHS
% ============================================================
\graphicspath{{figures/}{./figures/}{../figures/}}

% ============================================================
% SINGLE COLUMN FIGURE
% ============================================================
% Usage: \singlefig{filename}{caption}{label}

\newcommand{\singlefig}[3]{
\begin{figure}[htbp]
    \centering
    \includegraphics[width=\columnwidth]{#1}
    \caption{#2}
    \label{#3}
\end{figure}
}

% ============================================================
% DOUBLE COLUMN FIGURE
% ============================================================
% Usage: \doublefig{filename}{caption}{label}

\newcommand{\doublefig}[3]{
\begin{figure*}[htbp]
    \centering
    \includegraphics[width=\textwidth]{#1}
    \caption{#2}
    \label{#3}
\end{figure*}
}

% ============================================================
% SIDE-BY-SIDE FIGURES
% ============================================================
% Usage: \sidebyside{file1}{cap1}{lab1}{file2}{cap2}{lab2}{main_caption}{main_label}

\newcommand{\sidebyside}[8]{
\begin{figure}[htbp]
    \centering
    \begin{subfigure}[b]{0.48\columnwidth}
        \centering
        \includegraphics[width=\textwidth]{#1}
        \caption{#2}
        \label{#3}
    \end{subfigure}
    \hfill
    \begin{subfigure}[b]{0.48\columnwidth}
        \centering
        \includegraphics[width=\textwidth]{#4}
        \caption{#5}
        \label{#6}
    \end{subfigure}
    \caption{#7}
    \label{#8}
\end{figure}
}

% ============================================================
% THREE PANEL FIGURE
% ============================================================
% Usage: \threepanel{f1}{c1}{l1}{f2}{c2}{l2}{f3}{c3}{l3}{main_cap}{main_lab}

\newcommand{\threepanel}[2]{
\begin{figure*}[htbp]
    \centering
    #1
    \caption{#2}
\end{figure*}
}

% ============================================================
% HYPOTHESIS SUMMARY BAR CHART (TikZ)
% ============================================================
% Requires: \usepackage{pgfplots}
% Usage: Include data directly in document

% Example:
% \begin{figure}[htbp]
%     \centering
%     \begin{tikzpicture}
%         \begin{axis}[
%             ybar stacked,
%             bar width=15pt,
%             symbolic x coords={Domain1, Domain2, Domain3},
%             xtick=data,
%             xlabel={Research Domain},
%             ylabel={Number of Hypotheses},
%             legend style={at={(0.5,-0.15)}, anchor=north},
%         ]
%             \addplot[fill=shaktigreen] coordinates {(Domain1,8) (Domain2,9) (Domain3,7)};
%             \addplot[fill=shaktired] coordinates {(Domain1,2) (Domain2,3) (Domain3,3)};
%             \legend{Supported, Not Supported}
%         \end{axis}
%     \end{tikzpicture}
%     \caption{Summary of hypothesis test results by domain.}
%     \label{fig:hypothesis_summary}
% \end{figure}

% ============================================================
% BOX PLOT TEMPLATE
% ============================================================
% For efficiency comparisons - use with pgfplots

% Example:
% \begin{figure}[htbp]
%     \centering
%     \begin{tikzpicture}
%         \begin{axis}[
%             boxplot/draw direction=y,
%             xtick={1,2,3,4},
%             xticklabels={SHAKTI, Fixed, Uniform, CDA},
%             ylabel={Efficiency},
%         ]
%             \addplot+[boxplot prepared={
%                 median=0.94,
%                 upper quartile=0.96,
%                 lower quartile=0.91,
%                 upper whisker=0.99,
%                 lower whisker=0.85
%             }] coordinates {};
%             % Add more systems...
%         \end{axis}
%     \end{tikzpicture}
%     \caption{Efficiency comparison across trading mechanisms.}
%     \label{fig:efficiency_box}
% \end{figure}

% ============================================================
% SCALABILITY PLOT (LOG-LOG)
% ============================================================
% For showing O(n log n) scaling behavior

% Example:
% \begin{figure}[htbp]
%     \centering
%     \begin{tikzpicture}
%         \begin{loglogaxis}[
%             xlabel={Number of Agents},
%             ylabel={Clearing Time (ms)},
%             legend pos=north west,
%         ]
%             \addplot[only marks, mark=*] table {data/scalability.dat};
%             \addplot[domain=10:10000, samples=100, red, dashed] {x*ln(x)/100};
%             \legend{Observed, $O(n \log n)$}
%         \end{loglogaxis}
%     \end{tikzpicture}
%     \caption{Scalability analysis showing clearing time vs number of agents.}
%     \label{fig:scalability}
% \end{figure}

% ============================================================
% PARETO FRONT PLOT
% ============================================================
% For showing efficiency vs fairness trade-off

% Example:
% \begin{figure}[htbp]
%     \centering
%     \begin{tikzpicture}
%         \begin{axis}[
%             xlabel={Efficiency},
%             ylabel={Fairness (Gini Index)},
%             xmin=0.6, xmax=1.0,
%             ymin=0.6, ymax=1.0,
%         ]
%             \addplot[only marks, mark=o] coordinates {
%                 (0.78, 0.65) (0.82, 0.72) (0.89, 0.85) (0.94, 0.91)
%             };
%             \addplot[thick, shaktiblue, mark=*] coordinates {
%                 (0.94, 0.91)
%             };
%             \node[anchor=west] at (axis cs:0.94,0.91) {SHAKTI};
%         \end{axis}
%     \end{tikzpicture}
%     \caption{Pareto front showing efficiency-fairness trade-off.}
%     \label{fig:pareto}
% \end{figure}

% ============================================================
% VIOLIN PLOT APPROXIMATION
% ============================================================
% TikZ doesn't have native violin plots, use matplotlib export

% ============================================================
% FIGURE PLACEMENT HINTS
% ============================================================
% [h] - here (approximately)
% [t] - top of page
% [b] - bottom of page
% [p] - separate page for floats
% [!] - override internal parameters
%
% Recommended: [htbp] for flexibility

% ============================================================
% FIGURE SIZING GUIDELINES
% ============================================================
% Single column: \columnwidth (typically 3.5 inches)
% Double column: \textwidth (typically 7 inches)
% Specific sizes:
%   - Quarter page: 0.45\columnwidth
%   - Half page: 0.9\columnwidth
%   - Full width: \textwidth

% ============================================================
% RESOLUTION REQUIREMENTS
% ============================================================
% IEEE requires:
%   - Line art: 600 dpi minimum
%   - Halftones: 300 dpi minimum
%   - Color figures: 300 dpi minimum, CMYK color space
%
% Export settings:
%   - PDF for vector graphics (preferred)
%   - PNG at 300+ dpi for raster
%   - Avoid JPEG for figures with text/lines

% ============================================================
% ACCESSIBILITY
% ============================================================
% Ensure figures are interpretable in grayscale
% Use patterns in addition to colors
% Provide detailed captions describing key findings


% Required packages (should be in main document)
% \usepackage{graphicx}
% \usepackage{subcaption}
% \usepackage{tikz}
% \usepackage{pgfplots}

% ============================================================
% FIGURE PATHS
% ============================================================
\graphicspath{{figures/}{./figures/}{../figures/}}

% ============================================================
% SINGLE COLUMN FIGURE
% ============================================================
% Usage: \singlefig{filename}{caption}{label}

\newcommand{\singlefig}[3]{
\begin{figure}[htbp]
    \centering
    \includegraphics[width=\columnwidth]{#1}
    \caption{#2}
    \label{#3}
\end{figure}
}

% ============================================================
% DOUBLE COLUMN FIGURE
% ============================================================
% Usage: \doublefig{filename}{caption}{label}

\newcommand{\doublefig}[3]{
\begin{figure*}[htbp]
    \centering
    \includegraphics[width=\textwidth]{#1}
    \caption{#2}
    \label{#3}
\end{figure*}
}

% ============================================================
% SIDE-BY-SIDE FIGURES
% ============================================================
% Usage: \sidebyside{file1}{cap1}{lab1}{file2}{cap2}{lab2}{main_caption}{main_label}

\newcommand{\sidebyside}[8]{
\begin{figure}[htbp]
    \centering
    \begin{subfigure}[b]{0.48\columnwidth}
        \centering
        \includegraphics[width=\textwidth]{#1}
        \caption{#2}
        \label{#3}
    \end{subfigure}
    \hfill
    \begin{subfigure}[b]{0.48\columnwidth}
        \centering
        \includegraphics[width=\textwidth]{#4}
        \caption{#5}
        \label{#6}
    \end{subfigure}
    \caption{#7}
    \label{#8}
\end{figure}
}

% ============================================================
% THREE PANEL FIGURE
% ============================================================
% Usage: \threepanel{f1}{c1}{l1}{f2}{c2}{l2}{f3}{c3}{l3}{main_cap}{main_lab}

\newcommand{\threepanel}[2]{
\begin{figure*}[htbp]
    \centering
    #1
    \caption{#2}
\end{figure*}
}

% ============================================================
% HYPOTHESIS SUMMARY BAR CHART (TikZ)
% ============================================================
% Requires: \usepackage{pgfplots}
% Usage: Include data directly in document

% Example:
% \begin{figure}[htbp]
%     \centering
%     \begin{tikzpicture}
%         \begin{axis}[
%             ybar stacked,
%             bar width=15pt,
%             symbolic x coords={Domain1, Domain2, Domain3},
%             xtick=data,
%             xlabel={Research Domain},
%             ylabel={Number of Hypotheses},
%             legend style={at={(0.5,-0.15)}, anchor=north},
%         ]
%             \addplot[fill=shaktigreen] coordinates {(Domain1,8) (Domain2,9) (Domain3,7)};
%             \addplot[fill=shaktired] coordinates {(Domain1,2) (Domain2,3) (Domain3,3)};
%             \legend{Supported, Not Supported}
%         \end{axis}
%     \end{tikzpicture}
%     \caption{Summary of hypothesis test results by domain.}
%     \label{fig:hypothesis_summary}
% \end{figure}

% ============================================================
% BOX PLOT TEMPLATE
% ============================================================
% For efficiency comparisons - use with pgfplots

% Example:
% \begin{figure}[htbp]
%     \centering
%     \begin{tikzpicture}
%         \begin{axis}[
%             boxplot/draw direction=y,
%             xtick={1,2,3,4},
%             xticklabels={SHAKTI, Fixed, Uniform, CDA},
%             ylabel={Efficiency},
%         ]
%             \addplot+[boxplot prepared={
%                 median=0.94,
%                 upper quartile=0.96,
%                 lower quartile=0.91,
%                 upper whisker=0.99,
%                 lower whisker=0.85
%             }] coordinates {};
%             % Add more systems...
%         \end{axis}
%     \end{tikzpicture}
%     \caption{Efficiency comparison across trading mechanisms.}
%     \label{fig:efficiency_box}
% \end{figure}

% ============================================================
% SCALABILITY PLOT (LOG-LOG)
% ============================================================
% For showing O(n log n) scaling behavior

% Example:
% \begin{figure}[htbp]
%     \centering
%     \begin{tikzpicture}
%         \begin{loglogaxis}[
%             xlabel={Number of Agents},
%             ylabel={Clearing Time (ms)},
%             legend pos=north west,
%         ]
%             \addplot[only marks, mark=*] table {data/scalability.dat};
%             \addplot[domain=10:10000, samples=100, red, dashed] {x*ln(x)/100};
%             \legend{Observed, $O(n \log n)$}
%         \end{loglogaxis}
%     \end{tikzpicture}
%     \caption{Scalability analysis showing clearing time vs number of agents.}
%     \label{fig:scalability}
% \end{figure}

% ============================================================
% PARETO FRONT PLOT
% ============================================================
% For showing efficiency vs fairness trade-off

% Example:
% \begin{figure}[htbp]
%     \centering
%     \begin{tikzpicture}
%         \begin{axis}[
%             xlabel={Efficiency},
%             ylabel={Fairness (Gini Index)},
%             xmin=0.6, xmax=1.0,
%             ymin=0.6, ymax=1.0,
%         ]
%             \addplot[only marks, mark=o] coordinates {
%                 (0.78, 0.65) (0.82, 0.72) (0.89, 0.85) (0.94, 0.91)
%             };
%             \addplot[thick, shaktiblue, mark=*] coordinates {
%                 (0.94, 0.91)
%             };
%             \node[anchor=west] at (axis cs:0.94,0.91) {SHAKTI};
%         \end{axis}
%     \end{tikzpicture}
%     \caption{Pareto front showing efficiency-fairness trade-off.}
%     \label{fig:pareto}
% \end{figure}

% ============================================================
% VIOLIN PLOT APPROXIMATION
% ============================================================
% TikZ doesn't have native violin plots, use matplotlib export

% ============================================================
% FIGURE PLACEMENT HINTS
% ============================================================
% [h] - here (approximately)
% [t] - top of page
% [b] - bottom of page
% [p] - separate page for floats
% [!] - override internal parameters
%
% Recommended: [htbp] for flexibility

% ============================================================
% FIGURE SIZING GUIDELINES
% ============================================================
% Single column: \columnwidth (typically 3.5 inches)
% Double column: \textwidth (typically 7 inches)
% Specific sizes:
%   - Quarter page: 0.45\columnwidth
%   - Half page: 0.9\columnwidth
%   - Full width: \textwidth

% ============================================================
% RESOLUTION REQUIREMENTS
% ============================================================
% IEEE requires:
%   - Line art: 600 dpi minimum
%   - Halftones: 300 dpi minimum
%   - Color figures: 300 dpi minimum, CMYK color space
%
% Export settings:
%   - PDF for vector graphics (preferred)
%   - PNG at 300+ dpi for raster
%   - Avoid JPEG for figures with text/lines

% ============================================================
% ACCESSIBILITY
% ============================================================
% Ensure figures are interpretable in grayscale
% Use patterns in addition to colors
% Provide detailed captions describing key findings
